% GENERATED FILE. Edit canonical lessons, then run npm run generate:books.
% canonical-source-hash: fnv1a64:c0aee89ea43b3030
% canonical-lessons: UR-C519-dhol, UR-C519-saz, UR-C519-qavvali, UR-C519-mahfil, UR-C519-nachana

\chapter{Drums and Gatherings}
\label{ch:ur-drums-and-gatherings}
\hlchaptermodality{519}

\begin{chapteropening}
\textbf{By the end of this chapter:} \emph{I can talk about music at weddings and gatherings.}

Complete the last lesson of chapter 519: I can talk about music at weddings and gatherings.
\end{chapteropening}

\section[\texorpdfstring{\ur{ڈھول} (ḍhol)}{ڈھول (ḍhol)}]{\ur{ڈھول} (ḍhol) — a drum}

\label{lesson:UR-C519-dhol}



\begin{quote}
Before the new one: say the Urdu for practice (in music), then the Urdu for to hum.
\end{quote}

\subsection*{You'll want to know: \ur{ڈھول}}
\textbf{\ur{ڈھول}} — \emph{ḍhol} — \textquotedblleft{}a drum\textquotedblright{}.

\begin{quote}
\ur{شادی} \ur{میں} \ur{لڑکے} \ur{ڈھول} \ur{بجا} \ur{رہے} \ur{تھے۔}
\end{quote}

\begin{quote}
\emph{shādī meṅ laṛke ḍhol bajā rahe the.}
\end{quote}

\begin{quote}
At the wedding the boys were playing the drum.
\end{quote}

Masculine.

\begin{cousinweb}
\textbf{\ur{ڈھول}} (\emph{ḍhol}) is inherited from Sanskrit \emph{ḍhōla}, \textquotedblleft{}a large drum\textquotedblright{} (CDIAL 5608).
\end{cousinweb}

\begin{grammarlens}[title={What you've built}]
One more word for describing what you lived and felt.
\end{grammarlens}

\subsection*{Your turn}
\begin{itemize}
\raggedright
  \item \emph{Say it:} \emph{ḍhol}
  \item \emph{Say it:} \emph{ḍhol}, once more
  \item \emph{Say it:} \emph{shādī meṅ laṛke ḍhol bajā rahe the.} — At the wedding the boys were playing the drum.
  \item \emph{Recall:} say the Urdu for practice (in music), then the Urdu for to hum, then say \emph{ḍhol} again
\end{itemize}

\begin{tcolorbox}[breakable,colback=teal!4,colframe=teal!35!black,title={Before you move on}]
What does \ur{ڈھول} mean? (\textquotedblleft{}A drum\textquotedblright{}.) Say it once more.
\end{tcolorbox}
\section[\texorpdfstring{\ur{ساز} (sāz)}{ساز (sāz)}]{\ur{ساز} (sāz) — a musical instrument}

\label{lesson:UR-C519-saz}



\begin{quote}
Before the new one: say the Urdu for to hum, then the Urdu for a drum.
\end{quote}

\subsection*{You'll want to know: \ur{ساز}}
\textbf{\ur{ساز}} — \emph{sāz} — \textquotedblleft{}a musical instrument\textquotedblright{}.

\begin{quote}
\ur{میرا} \ur{بھائی} \ur{تین} \ur{ساز} \ur{بجاتا} \ur{ہے۔}
\end{quote}

\begin{quote}
\emph{merā bhāī tīn sāz bajātā hai.}
\end{quote}

\begin{quote}
My brother plays three instruments.
\end{quote}

Masculine.

\begin{cousinweb}
\textbf{\ur{ساز}} (\emph{sāz}) is usually traced to Persian \emph{sāz}, \textquotedblleft{}equipment; a musical instrument\textquotedblright{}.
\end{cousinweb}

\begin{grammarlens}[title={What you've built}]
One more word for describing what you lived and felt.
\end{grammarlens}

\subsection*{Your turn}
\begin{itemize}
\raggedright
  \item \emph{Say it:} \emph{sāz}
  \item \emph{Say it:} \emph{sāz}, once more
  \item \emph{Say it:} \emph{merā bhāī tīn sāz bajātā hai.} — My brother plays three instruments.
  \item \emph{Recall:} say the Urdu for to hum, then the Urdu for a drum, then say \emph{sāz} again
\end{itemize}

\begin{tcolorbox}[breakable,colback=teal!4,colframe=teal!35!black,title={Before you move on}]
What does \ur{ساز} mean? (\textquotedblleft{}A musical instrument\textquotedblright{}.) Say it once more.
\end{tcolorbox}
\section[\texorpdfstring{\ur{قوالی} (qavvālī)}{قوالی (qavvālī)}]{\ur{قوالی} (qavvālī) — qawwali (devotional song)}

\label{lesson:UR-C519-qavvali}



\begin{quote}
Before the new one: say the Urdu for a drum, then the Urdu for a musical instrument.
\end{quote}

\subsection*{You'll want to know: \ur{قوالی}}
\textbf{\ur{قوالی}} — \emph{qavvālī} — \textquotedblleft{}qawwali (devotional song)\textquotedblright{}.

\begin{quote}
\ur{ہم} \ur{نے} \ur{کل} \ur{رات} \ur{قوالی} \ur{سنی۔}
\end{quote}

\begin{quote}
\emph{ham ne kal rāt qavvālī sunī.}
\end{quote}

\begin{quote}
Last night we listened to qawwali.
\end{quote}

Feminine.

\begin{cousinweb}
\textbf{\ur{قوالی}} (\emph{qavvālī}) is built on \emph{qavvāl}, \textquotedblleft{}a qawwali singer\textquotedblright{}, usually traced, through Persian, to Arabic \emph{qawwāl}, \textquotedblleft{}one who says much\textquotedblright{}, with the noun ending \emph{-ī}.
\end{cousinweb}

\begin{grammarlens}[title={What you've built}]
One more word for describing what you lived and felt.
\end{grammarlens}

\subsection*{Your turn}
\begin{itemize}
\raggedright
  \item \emph{Say it:} \emph{qavvālī}
  \item \emph{Say it:} \emph{qavvālī}, once more
  \item \emph{Say it:} \emph{ham ne kal rāt qavvālī sunī.} — Last night we listened to qawwali.
  \item \emph{Recall:} say the Urdu for a drum, then the Urdu for a musical instrument, then say \emph{qavvālī} again
\end{itemize}

\begin{tcolorbox}[breakable,colback=teal!4,colframe=teal!35!black,title={Before you move on}]
What does \ur{قوالی} mean? (\textquotedblleft{}Qawwali (devotional song)\textquotedblright{}.) Say it once more.
\end{tcolorbox}
\section[\texorpdfstring{\ur{محفل} (mahfil)}{محفل (mahfil)}]{\ur{محفل} (mahfil) — a gathering (for music or poetry)}

\label{lesson:UR-C519-mahfil}



\begin{quote}
Before the new one: say the Urdu for a musical instrument, then the Urdu for qawwali.
\end{quote}

\subsection*{You'll want to know: \ur{محفل}}
\textbf{\ur{محفل}} — \emph{mahfil} — \textquotedblleft{}a gathering (for music or poetry)\textquotedblright{}.

\begin{quote}
\ur{کل} \ur{رات} \ur{دادا} \ur{کے} \ur{گھر} \ur{محفل} \ur{تھی۔}
\end{quote}

\begin{quote}
\emph{kal rāt dādā ke ghar mahfil thī.}
\end{quote}

\begin{quote}
Last night there was a musical evening at grandfather's house.
\end{quote}

Feminine.

\begin{cousinweb}
\textbf{\ur{محفل}} (\emph{mahfil}) is usually traced, through Persian, to Arabic \emph{maḥfil}, \textquotedblleft{}a place where people gather\textquotedblright{}.
\end{cousinweb}

\begin{grammarlens}[title={What you've built}]
One more word for describing what you lived and felt.
\end{grammarlens}

\subsection*{Your turn}
\begin{itemize}
\raggedright
  \item \emph{Say it:} \emph{mahfil}
  \item \emph{Say it:} \emph{mahfil}, once more
  \item \emph{Say it:} \emph{kal rāt dādā ke ghar mahfil thī.} — Last night there was a musical evening at grandfather's house.
  \item \emph{Recall:} say the Urdu for a musical instrument, then the Urdu for qawwali, then say \emph{mahfil} again
\end{itemize}

\begin{tcolorbox}[breakable,colback=teal!4,colframe=teal!35!black,title={Before you move on}]
What does \ur{محفل} mean? (\textquotedblleft{}A gathering (for music or poetry)\textquotedblright{}.) Say it once more.
\end{tcolorbox}
\section[\texorpdfstring{\ur{نچانا} (nachānā)}{نچانا (nachānā)}]{\ur{نچانا} (nachānā) — to make (someone) dance}

\label{lesson:UR-C519-nachana}



\begin{quote}
Before the new one: say the Urdu for qawwali, then the Urdu for a gathering (for music or poetry).
\end{quote}

\subsection*{You'll want to know: \ur{نچانا}}
\textbf{\ur{نچانا}} — \emph{nachānā} — \textquotedblleft{}to make (someone) dance\textquotedblright{}.

\begin{quote}
\ur{ڈھول} \ur{نے} \ur{سب} \ur{مہمانوں} \ur{کو} \ur{نچایا۔}
\end{quote}

\begin{quote}
\emph{ḍhol ne sab mehmānoṅ ko nachāyā.}
\end{quote}

\begin{quote}
The drum got all the guests dancing.
\end{quote}

Past, with \ur{نے} (he or she …): \ur{اس} \ur{نے} \ur{بچوں} \ur{کو} \ur{نچایا}, \emph{us ne bachchoṅ ko nachāyā}.

\begin{cousinweb}
\textbf{\ur{نچانا}} (\emph{nachānā}) is the causative of \emph{nāchnā}, \textquotedblleft{}to dance\textquotedblright{}, which CDIAL lists under Sanskrit \emph{nṛtyati}, \textquotedblleft{}dances\textquotedblright{} (CDIAL 7583).
\end{cousinweb}

\begin{grammarlens}[title={What you've built}]
One more word for describing what you lived and felt.
\end{grammarlens}

\subsection*{Your turn}
\begin{itemize}
\raggedright
  \item \emph{Say it:} \emph{nachānā}
  \item \emph{Say it:} \emph{nachānā}, once more
  \item \emph{Say it:} \emph{ḍhol ne sab mehmānoṅ ko nachāyā.} — The drum got all the guests dancing.
  \item \emph{Recall:} say the Urdu for qawwali, then the Urdu for a gathering (for music or poetry), then say \emph{nachānā} again
\end{itemize}

\begin{tcolorbox}[breakable,colback=teal!4,colframe=teal!35!black,title={Before you move on}]
What does \ur{نچانا} mean? (\textquotedblleft{}To make (someone) dance\textquotedblright{}.) Say it once more.
\end{tcolorbox}
